\section{Shared-input convex vector graphs}\label{sec:vector}

Several outputs sharing an input must be approximated at the same input
point. Separate scalar formulations need not achieve a good joint integer
count. We first construct common partitions, then replace output count by
the dimension of the component curvatures. All finite statements allow
real coefficients and unrestricted finite continuous size. All computational
statements in this section assume dense rational polynomial data and
positive rational tolerances. They concern polynomial construction time
and total rational encoding, including the supplied error-body data.

For a vector function on an interval, write $g_F=T_F-F$ for its vector
chord gap. Componentwise convexity means $g_F\ge0$ coordinatewise.
For continuous functions, curvature rank means linear dimension modulo
affine functions; no derivative is required. The first comparison pays for
the number of outputs. Selecting a basis from the original convex outputs
replaces that dependence by curvature rank. For general unconditional
bodies, positive polar directions control the error and an explicit inner
polytope supplies the final linear band. A shared basis then extends the
construction to separable inputs. The separation examples distinguish
which hypotheses each comparison needs.

\subsection{Level refinement and an implicit common partition}

\begin{lemma}[Refinement at an arbitrary factor]\label{lem:level-refinement}
If a continuous convex function has chord error at most $H\tau$ on an
interval, where $H\ge1$ is an integer and $\tau>0$, it has a partition
into at most $2H-1$ intervals with chord error at most $\tau$.
An interval cover by $N$ intervals having a given scalar or componentwise
convex chord-error bound can be replaced by a partition with at most
$N$ intervals and the same bound.
\end{lemma}
\begin{proof}
Let $g$ be the original nonnegative concave chord gap. For each level
$j\tau$, $1\le j<H$, strictly below its maximum, cut at both endpoints
of $\{g\ge j\tau\}$. Along the two sides of the maximum, consecutive
cuts enclose a range of $g$ of width at most $\tau$. The middle interval
has equal endpoint values at the highest retained level and maximum
at most one further level. Its range also has width at most $\tau$.
If no level is retained the original maximum is at most $\tau$.
On a new interval the local chord gap equals $g$ minus its endpoint
interpolant, hence is nonnegative and at most this range width.
There are at most $2H-2$ cuts. Plateaus, repeated cuts and singleton
intervals cause no difficulty.

For the cover statement, begin at the left endpoint and select a covering
interval containing that point and extending farthest right. At each new
endpoint repeat this choice among covering intervals containing it.
Unless the right domain endpoint has been reached, the finite covering
property ensures strict progress. No selected interval is needed twice.
Restrict each selected interval to the new partition cell. Convex chord
errors decrease under restriction, component by component.
\end{proof}

\begin{proposition}[Finite output-count comparison]\label{prop:finite-vector-overlay}
For continuous componentwise convex $F:[0,1]\to\R^m$ and positive box
tolerances,
\[
 \pbin\le\pconv+\lceil\log_2(2m+1)\rceil.
\]
For $K=\{e:A|e|\le b\}$ with $A\ge0$, $b>0$, and $q$ rows, compactness
of $K$ gives the finite bound
$\pbin\le\pconv+\lceil\log_2(8q+1)\rceil$.
\end{proposition}
\begin{proof}
On the interval hull of a parity contact from \cref{lem:parity},
each component has midpoint gap at most $\eps_j$, and hence full chord
gap at most $2\eps_j$. The three-piece proof in
\cref{lem:scalar-chords} uses at most two cuts for each component.
Overlaying these cuts yields at most $2m+1$ cells, all with
$0\le g_j\le\eps_j$. The bands
$T_j-\eps_j\le w_j\le T_j$ contain the exact vector graph and admit
only the prescribed errors. At most $(2m+1)2^p$ bands result from a
$p$-integer lift; \cref{lem:disjunction} proves the first bound.
For a facet row use the convex function $(AF)_k/b_k$. Its gap on a
parity hull is at most two. Two three-piece refinements give gap at
most $1/2$ using at most eight cuts. Overlaying all rows gives at most
$8q+1$ cells with $g_F\in K/2$. The band $w-T_F\in K/2$ contains
the graph and admits errors in $K/2+K/2=K$. It is linear using
continuous absolute-value auxiliaries. Degenerate hulls can be kept
as singleton cells throughout.
\end{proof}

The half-body band matters for coupled budgets. A downward shift by an
arbitrary nonnegative vector in $K$ can produce an error outside $K$:
$(1,0)$ and $(0,1)$ both belong to the unit $\ell_1$ ball, but their
difference has norm two.

\begin{lemma}[Sorted random-access overlay]\label{lem:implicit-overlay}
Suppose polynomially many nondecreasing rational arrays
$0=a_{j,0}\le\cdots\le a_{j,K_j}=1$ have polynomial-time indexed
evaluators. Assume $K_j$ have polynomial bit length and all knots are
integer multiples of $1/H$ for a known positive integer $H$ of polynomial
bit length. The sorted multiset of all right endpoints
$a_{j,1},\ldots,a_{j,K_j}$, with a single zero prepended, has a
polynomial-time indexed evaluator. Its $K=\sum_jK_j$ cells cover the
interval, and every positive-length cell lies inside a source cell
of every array. Duplicates are retained as zero-length cells.
\end{lemma}
\begin{proof}
For an integer $v\in[0,H]$, binary search each source index to count
its right endpoints at most $v/H$. Sum these counts. The numerator
of the $k$th multiset endpoint is the least $v$ whose count is at
least $k$, found by binary search in $[0,H]$. This uses
$O(\log(H+1)\sum_j\log(K_j+1))$ indexed evaluations and polynomial
rational arithmetic. Multiplicity makes equality an exact integer
comparison. If a source knot lay strictly inside a positive merged
cell it would also appear in the merged multiset, a contradiction.
For source-cell metadata at a left endpoint $a<1$, take the last source
knot at most $a$ and its successor; the successor is at least the right
merged endpoint. At a repeated merged endpoint any source cell
containing that point suffices; at one use the final source cell.
\end{proof}

The mass-knot procedure in \cref{thm:compiled-curvature} can be made
nondecreasing without increasing its precision. Use the same deterministic
mass evaluation and accuracy at each bisection node for every target.
At a node the targets going left, returning the midpoint, and going right
are ordered ranges. All outputs of its left subtree are at most the
midpoint, and all outputs of its right subtree are at least it. Induction
on the fixed tree depth proves monotonicity, even if the approximate mass
values at different nodes are not monotone. Forced endpoints preserve
order. Reverse local indices when reflecting a piece. Greedy knots are
already increasing; concatenating the hybrid's consecutive pieces
therefore gives the arrays required in \cref{lem:implicit-overlay}.
The common denominator is the product of the polynomially many rational
piece-endpoint denominators times the largest local dyadic denominator.
Its bit length is polynomial; individual short denominators alone would
not establish this common-grid property.

\begin{theorem}[Compiled convex vector overlay]\label{thm:convex-vector-overlay}
For dense rational componentwise convex polynomial outputs with positive
rational box tolerances, a polynomial rational MILP has
\[
 \pout\le\pconv+11+\lceil\log_2m\rceil.
\]
An entirely affine graph has an exact zero-integer formulation.
\end{theorem}
\begin{proof}
Use the sorted scalar hybrid arrays at tolerance $\eps_j$. Their actual
counts obey $K_j\le1458\,2^p$ by \eqref{eq:hybrid-total-cells}, because
projecting any $p$-integer vector lift gives a scalar lift. Each source
cell has chord gap at most $13\eps_j/16$. An affine component can use
one cell. Apply \cref{lem:implicit-overlay} to their right endpoints.
Evaluate each original polynomial at each selected merged endpoint
and round downward with error at most $\eps_j/8$, using fixed dyadic
output denominators and offsets for signed values. Let $y_j$ interpolate
these rounded values at the common continuous weight. Since restriction
preserves convex chord error, the band
\begin{equation}\label{eq:vector-directed-band}
 y_j-13\eps_j/16\le w_j\le y_j+\eps_j/8
\end{equation}
contains the graph and admits absolute error at most $15\eps_j/16$.
This remains true at zero-length cells. All components use the same
input and interpolation weight. \Cref{lem:indexed-compiler} needs
only $\lceil\log_2K\rceil$ binary index variables: the endpoint
arithmetic and internal Boolean wires remain continuous. Exclude codes
at least $K$. Since $K\le1458m2^p<2048m2^p$, the count follows.
The algorithm computes actual counts and never uses the comparator lift.
\end{proof}

Merging explicit breakpoints and sharing an SOS2 index for several outputs
are established formulation tools; \citet[Section 3, Proposition 1]{lyu2026}
give that common-input construction. The point of the present overlay is
its polynomial-bit random access when the underlying lists can be
exponentially long, together with a comparison against every convex
integer lift. No sparse huge-degree evaluation claim is made here.

\subsection{Original-output curvature bases}

\begin{lemma}[Barycentric spanners of curvature rows]\label{lem:curvature-spanner}
Let $G_1,\ldots,G_m$ be continuous convex functions whose classes modulo
affine functions span dimension $r>0$. There are $r$ original functions
$G_{i_s}$ such that
\begin{equation}\label{eq:curvature-basis}
 G_j=\alpha_j+\beta_jx+\sum_{s=1}^r c_{js}G_{i_s},
 \qquad |c_{js}|\le1.
\end{equation}
For rational dense polynomials, a polynomial-time rational algorithm
returns original functions with the weaker bound $|c_{js}|\le2$.
If $\Psi=\sum_sG_{i_s}$ and the coefficient bound is $c$, then every
interval satisfies $0\le g_{G_j}\le c g_\Psi$.
\end{lemma}
\begin{proof}
Choose coordinates on the quotient and maximize the absolute determinant
of $r$ original rows. Replacing one row by any other multiplies the
determinant by its representation coefficient, so Cramer's rule proves
\eqref{eq:curvature-basis}. Affine terms cancel in chord gaps. The selected
gaps are nonnegative, hence signed coefficients of magnitude at most $c$
still give $g_{G_j}\le c\sum_sg_{G_{i_s}}$.

For rational polynomials form the nonlinear coefficient matrix and
retain $r$ independent columns. Start with a nonsingular row minor.
Whenever a representation coefficient has magnitude greater than two,
replace the corresponding basis row. Its determinant more than doubles.
If the retained matrix has total encoding length $L$, the product of
its entry denominators is at most $2^L$. Nonzero row determinants are
at least $2^{-rL}$ and all are at most $r!2^{rL}$. Thus at most
$2rL+\log_2(r!)+1$ exchanges occur. Every current basis is a submatrix
of the original matrix; inverses and exact comparisons have polynomial
bit length. Relations on independent columns extend to all columns,
so the resulting coefficients represent the entire polynomial rows.
\end{proof}

The maximum-determinant argument and determinant exchange are the
barycentric-spanner constructions of \citet[Section 2.3,
Propositions 2.2 and 2.4]{awerbuch2004}. The requirement to select original
convex outputs, rather than an arbitrary algebraic basis, preserves
nonnegative chord gaps.

\begin{theorem}[Box and facet curvature rank]\label{thm:vector-rank}
For box tolerances normalize $G_j=F_j/\eps_j$, and let $r$ be the
dimension of their span modulo affine functions. For $r\ge1$,
\begin{equation}\label{eq:vector-box-rank}
 \pbin\le\pconv+\lceil\log_2(4r-1)\rceil,
 \qquad
 \pout\le\pconv+\lceil\log_2r\rceil+12.
\end{equation}
The first bound is finite for continuous convex components; the second
is a polynomial rational construction for dense polynomial data.
For a compact body $K=\{e:A|e|\le b\}$, $A\ge0$, $b>0$, let $r$ instead
be the curvature rank of $H_k=(AF)_k/b_k$. Then
\begin{equation}\label{eq:vector-facet-rank}
 \pbin\le\pconv+\lceil\log_2(8r-1)\rceil,
 \qquad
 \pout\le\pconv+\lceil\log_2r\rceil+13.
\end{equation}
The computational facet statement includes rational $A,b$ in its input.
In either model $r=0$ forces an affine graph and needs no integers.
\end{theorem}
\begin{proof}
Choose an exact original-output basis and sum it to obtain $\Psi$.
Each normalized basis function has midpoint gap at most one on a
parity hull, so $g_\Psi\le2r$ there. For a box, refine to scalar gap
one using at most $4r-1$ cells. \Cref{lem:curvature-spanner} bounds
all normalized component gaps by one, and downward chord bands give
the finite claim. For facets refine to $1/2$, using at most $8r-1$
cells. Then $g_F\in K/2$; the symmetric band $w-T_F\in K/2$ gives
containment and admissibility.

For a rational box use the factor-two basis, so
$g_{G_j}\le2g_\Psi$. Refining parity hulls to tolerance $1/2$ gives
$N_{1/2}(\Psi)\le(8r-1)2^p$ by \cref{lem:level-refinement}.
The hybrid has at most $486N_{1/2}$ cells, each with
$g_\Psi\le13/32$. Evaluate all original normalized outputs at their
common endpoints, round downward within $1/8$, and use
\eqref{eq:vector-directed-band} with unit normalized tolerances.
The cell count is less than $3888r2^p<2^{12}r2^p$.

For rational facets use the factor-two basis of the $H_k$ and tolerance
$1/4$. Now $N_{1/4}(\Psi)\le(16r-1)2^p$ and every hybrid cell has
$g_\Psi\le13/64$, hence $g_F\in(13/32)K$. Put
\[
 M_A=\max_k\sum_j A_{kj}/b_k>0,
 \qquad \delta=\min\{1,(16M_A)^{-1}\}.
\]
Round each original endpoint coordinate downward within $\delta$.
Its rounding error belongs to $K/16$, as does the interpolated error.
The resulting center $y$ satisfies $y-F\in(15/32)K$.
The band $w-y\in K/2$ contains the graph and admits errors in
$(31/32)K$. Its explicit linear form is
\begin{equation}\label{eq:facet-linear-band}
 s_j\ge w_j-y_j,\quad s_j\ge y_j-w_j,\quad As\le b/2,
 \qquad s\text{ continuous}.
\end{equation}
The count is less than $486\cdot16r2^p<2^{13}r2^p$.
All additional component evaluations and bands add only continuous
variables to the shared index compiler.

Finally, rank zero in the box model means each output is affine.
For facets, compactness and nonnegativity imply that each column of
$A$ has a positive entry. Vanishing nonnegative weighted chord gaps
therefore force every component gap to vanish. This also proves the
rank-zero assertion for continuous functions.
\end{proof}

For weighted $\ell_1$ error, $K=\{e:\sum_j a_j|e_j|\le b\}$ with
$a_j>0$, there is one facet scalarization. Unless the graph is affine,
its rank is one; \eqref{eq:vector-facet-rank} has finite overhead three
and compiled overhead thirteen, independent of output count and degree.
These are upper comparisons, not assertions that the logarithmic rank
loss is necessary.

\subsection{Direct allocation for unconditional bodies}

\begin{proposition}[Maximum-product box transfer]\label{prop:vector-log-product}
For any unconditional error body $K\subset\R^m$ and continuous
componentwise convex curve,
\[
 \pbin\le\pconv+\lceil\log_2(4m-1)\rceil.
\]
For dense rational polynomial components and a polynomial rational
strong separation oracle for $K$, with known positive rational radii
$\rho B_2^m\subset K\subset R B_\infty^m$, a polynomial rational
construction has $\pout\le\pconv+14+\lceil\log_2m\rceil$.
The final formulation uses explicit linear bands.
\end{proposition}
\begin{proof}
Maximize $\prod_j b_j$ over nonnegative $b\in K$. A maximizer is
strictly positive, because $K$ is compact and contains a positive
point. Sign changes and convexity give $[-b,b]\subset K$.
Differentiating the log product along a feasible segment from $b$ to
$v\in K\cap\R_+^m$ yields
\begin{equation}\label{eq:product-support}
 \sum_jv_j/b_j\le m.
\end{equation}
Thus $\Psi=\sum_jF_j/b_j$ has full gap at most $2m$ on every parity
hull. Refine it to gap one with at most $4m-1$ cells. Every normalized
component gap is at most one, so downward bands in $[-b,b]$ prove
the finite statement.

The scalar specialization of \cref{lem:block-logdet-oracle}, applied
to $RI_m$ and coordinate cap one, returns rational $b>0$ in $K$ with
product at least $e^{-1}$ times optimal. The cap loses no positive
feasible point, since all coordinates of $K$ are at most $R$.
This approximate maximizer obeys
\begin{equation}\label{eq:approximate-product-support}
 \sum_jv_j/b_j\le7m\qquad(v\in K\cap\R_+^m).
\end{equation}
Indeed for $m\ge2$ set $s=1-1/m$, $a_j=v_j/b_j$, $S=\sum_ja_j$.
The feasible point $sb+v/m$ has product ratio at most $e$, whereas
expansion with nonnegative terms gives ratio at least
$s^{m-1}(s+S/m)$. Since $s^{m-1}\ge e^{-1}$,
$S\le m(e^2-1)+1<7m$. For $m=1$ the ratio bound directly gives
$S\le e<7$.

For this rational $\Psi$ the parity-hull gap is at most $14m$;
therefore $N_1(\Psi)\le(28m-1)2^p$. Its hybrid cells have scalar gap
at most $13/16$. Each component $F_j/b_j$ has no larger gap.
Evaluate and round these normalized values downward within $1/8$
and impose the normalized version of \eqref{eq:vector-directed-band}.
The error belongs to $(15/16)[-b,b]\subset K$ and the graph is
contained. The count is below $486\cdot28m2^p<2^{14}m2^p$.
The rational oracle bounds $b$, $1/b_j$, and all subsequent precision
lengths polynomially. An affine graph is retained exactly.
\end{proof}

Maximum-product allocation and the first-order inequality
\eqref{eq:product-support} are classical; in network allocation they
are the proportional-fairness relation of \citet[Section 2]{kelly1998}.
This direct construction avoids any facet enumeration. The next result
uses the effective nonlinear image to replace its output-count penalty.

\section{Curvature rank with a strong error-body oracle}\label{sec:vector-oracle}

Let $K$ be unconditional and put
\begin{equation}\label{eq:nonlinear-output-image}
 S=\operatorname{span}\{F(x)-(1-x)F(0)-xF(1):0\le x\le1\},
 \qquad r=\dim S.
\end{equation}
Choose a full-column-rank $m\times r$ matrix $V$ spanning $S$ and write
$F(x)=a+bx+Vq(x)$, with $q(0)=q(1)=0$.
For rational dense polynomials, select independent nonlinear coefficient
columns for $V$ and use a rational left inverse $J$ to obtain $q$.
Elimination has polynomial bit cost. The effective body
$C=\{z:Vz\in K\}$ is symmetric, compact, and full dimensional in
$\R^r$, although it need not be unconditional there.
If $\rho B_2^m\subset K\subset R B_\infty^m$, known radii are
\begin{equation}\label{eq:effective-body-radii}
 \frac{\rho}{1+\sum|V_{ij}|}B_2^r\subset C
 \subset mR(1+\sum|J_{ij}|)B_2^r.
\end{equation}
A separator for $K$ at $Vz$ pulls back by $V^T$. At an exterior query
its pulled-back normal is nonzero, because zero belongs to $C$.

\subsection{Exactly feasible rational spanners}

We give the oracle interface in detail, because approximate optimization
alone neither proves exact feasibility nor controls denominator growth
through repeated basis exchanges. The imported weak-optimization theorem
is \citet[Definition (5), Theorem (3.1)]{gls1981}: given a rational weak
separator and explicit inner and outer balls for a convex body $P$, it
returns a rational $z$ with
\begin{equation}\label{eq:weak-opt-spanner}
 \operatorname{dist}(z,P)\le\eta,
 \qquad d^Tz\ge\max_{x\in P}d^Tx-\eta.
\end{equation}
The objective comparison is with all of $P$. In dimension one apply the
theorem to $P\times[-1,1]$ and objective $(d,0)$, then project; inner
radius $\min(\sigma,1)$ and outer radius $R_0+1$ suffice.

\begin{lemma}[Fixed-grid feasible spanner]\label{lem:rational-spanner}
Let $B_2(c,\sigma)\subset P\subset B_2(c,R_0)$ be a compact convex
body with a polynomial rational weak separator and explicit rational
$c,\sigma,R_0$, $0<\sigma\le R_0$. Let $W\in\Q^{d\times r}$,
$r\ge1$, and suppose rational seeds $\lambda_1,\ldots,\lambda_r\in P$
have independent image rows $\lambda_i^TW$. A polynomial-time rational
algorithm returns feasible such rows which represent every
$\lambda^TW$, $\lambda\in P$, with coefficients of magnitude at most
$9/4<3$. Complexity is polynomial in the entire input, balls, and seeds.
\end{lemma}
\begin{proof}
Let $B_0$ be the seed row matrix and $\Delta_0=|\det B_0|>0$.
Choose the rational bounds
\[
 Z=1+\sum_i|c_i|+R_0,\quad M=1+\sum_{ij}|W_{ij}|,\quad
 A=1+dMZ,\quad U=1+r!A^{r-1}/\Delta_0,\quad L=1+dMU.
\]
Every feasible image coordinate has magnitude at most $A$.
As long as determinants never decrease, cofactors bound every inverse
basis entry by $U$. Each representation-coordinate objective
$d_i=WB^{-1}e_i$ and its negative then has $\ell_1$ norm at most $L$.
These bounds have polynomial bit length, including when $r$ varies.

Choose one positive dyadic $\delta$ before starting exchanges, such that
\begin{equation}\label{eq:fixed-spanner-grid}
 2d\delta\le1,\qquad
 d\delta[1+L+2L(R_0+1)/\sigma]\le1/4.
\end{equation}
For any objective $d'$ of $\ell_1$ norm at most $L$, apply
\eqref{eq:weak-opt-spanner} with $\eta=d\delta$. Round its output
coordinatewise to multiples of $\delta$, obtaining $q$ with
$\|q-z\|_2\le d\delta$, and set
\begin{equation}\label{eq:spanner-repair}
 \epsilon=2d\delta,\qquad
 \lambda=\frac{\sigma q+\epsilon c}{\sigma+\epsilon}.
\end{equation}
Since $\operatorname{dist}(q,P)\le\epsilon$, choose $x\in P$ within
that distance. The repaired point is the convex combination, with
weights $\sigma/(\sigma+\epsilon)$ and
$\epsilon/(\sigma+\epsilon)$, of $x$ and
$c+(\sigma/\epsilon)(q-x)\in B_2(c,\sigma)$.
Thus $\lambda\in P$ exactly. Also $\|c-q\|_2\le R_0+1$, so
\[
 \max_{x\in P}(d')^Tx-(d')^T\lambda
 \le d\delta+L[d\delta+2d\delta(R_0+1)/\sigma]\le1/4.
\]
Every returned point uses the same grid and the same repair weights.
Its coordinates have a fixed common rational denominator of polynomial
bit length and uniformly bounded numerators. Including seed denominators
preserves this statement for every possible basis.

At each basis position optimize both signs of $d_i$. If any returned
coefficient has magnitude greater than two, exchange that row. The
determinant more than doubles. Otherwise every exact feasible coefficient
has magnitude at most $2+1/4$, proving the spanner guarantee.
Since all determinants are at most $r!A^r$, there are at most
$1+\lceil\log_2(r!A^r/\Delta_0)\rceil$ exchanges. Fixed denominators
ensure that every inverse, later objective, oracle input, and comparison
retains polynomial bit length. This proves uniform polynomial complexity,
not merely a per-call bound in a potentially growing current input.
\end{proof}

To apply the lemma to positive polar directions, a weak separator must
itself be constructed. Replace the coordinate outer radius by the rational
Euclidean radius $R_E=mR$, so
$\rho B_2^m\subset K\subset R_EB_2^m$.
For rational $a$ and $\tau>0$, weak optimization over $K$ followed by
central repair returns rational $x\in K$ with
\begin{equation}\label{eq:feasible-support}
 a^Tx\le h_K(a)\le a^Tx+\tau.
\end{equation}
Indeed take $\eta\le1$ satisfying
$\eta[1+\|a\|_1(R_E+1)/\rho]\le\tau$ in
\eqref{eq:weak-opt-spanner} and put $x=\rho z/(\rho+\eta)$.
The inner-ball convex-combination argument proves exact feasibility;
the support loss is at most the displayed quantity. The zero objective
can return zero.

Define the full-dimensional positive polar
\[
 P_+=\{\lambda\ge0:h_K(\lambda)\le1\}.
\]
Choose a positive dyadic $a_0\le\min\{1,(8R_Em)^{-1}\}$.
An inner ball has center $c=a_0\mathbf1$ and radius $a_0/2$:
its coordinates are positive and its Euclidean norm is less than
$1/R_E$. An outer ball centered at $c$ has radius
$1/\rho+ma_0$, since $P_+\subset\rho^{-1}B_2^m$.
For a rational query $\lambda$ and distance accuracy $\eta>0$, first
separate any coordinate outside $[0,1/\rho]$ by the appropriate
coordinate inequality. Otherwise set
$\tau=\eta/[2(1+m/\rho)]$ and obtain $x$ from
\eqref{eq:feasible-support}. If $\lambda^Tx>1$, the inequality
$x^T\mu\le1$ separates the query from all of $P_+$.
Its rational nonzero normal can be rescaled by its infinity norm to
meet the norm convention in the weak separator. If
$\lambda^Tx\le1$, then $h_K(\lambda)\le1+\tau$ and
$\lambda/(1+\tau)\in P_+$ lies within
$\tau\|\lambda\|_2\le\eta/2$ of the query.
This is a valid weak-membership answer, not an exact membership test.
All accuracies and returned coefficients have polynomial bit length.

Choose independent rows $j_1,\ldots,j_r$ of $V$.
The feasible seeds $2a_0e_{j_s}\in P_+$ have independent image rows
$2a_0e_{j_s}^TV$. \Cref{lem:rational-spanner} therefore constructs
feasible nonnegative $\lambda_s$ whose images span and
$3$-approximately represent all $\lambda^TV$, $\lambda\in P_+$.
For $C$ in \eqref{eq:effective-body-radii}, use its strong separator,
center zero, image matrix $I_r$, and small coordinate-axis seeds from
its inner ball. The same lemma returns a feasible factor-three basis
of $C$. This derives both oracle calls from the original strong oracle;
it assumes neither a polyhedral body nor a strong separator for its polar.
Positive-polar access is a quantitative implementation of the classical
polar and nonnegative anti-blocker oracle equivalences
\citep[Corollaries (3.4)--(3.5)]{gls1981}. Approximate optimization
combined with barycentric-spanner exchange also predates this construction;
\citet[Section 3.3]{plevrakis2020} use that combination in online control.
The explicit fixed denominator and exact-feasibility repair here supply
the rational interface needed by the graph compiler, without asserting
a new general spanner or polar-oracle principle.

\subsection{A linear band in the effective image}

The two spanners serve different purposes. The positive polar spanner
reduces all output budgets to one convex scalar function. The spanner of
the effective body turns that scalar gap certificate into an explicit
linear band. Summing the polar directions contributes one rank factor to
the refinement count, and the inner polytope contributes another. Together
they yield an overhead proportional to twice the logarithm of the rank.
The band's auxiliary coordinates remain
continuous, so the approximation index is still the only source of binaries.

\begin{theorem}[Oracle-body curvature rank]\label{thm:vector-oracle-rank}
For a continuous componentwise convex curve and unconditional error body,
with $r$ as in \eqref{eq:nonlinear-output-image} and $r\ge1$,
\begin{equation}\label{eq:oracle-rank-count}
 \pbin\le\pconv+\lceil\log_2(8r^2-1)\rceil.
\end{equation}
Under the dense rational and strong-oracle hypotheses above, a polynomial
rational MILP has
\begin{equation}\label{eq:oracle-rank-compiled}
 \pout\le\pconv+17+2\lceil\log_2r\rceil.
\end{equation}
No oracle constraint remains in that MILP. For $r=0$ the graph is affine.
\end{theorem}
\begin{proof}
For nonnegative $e$, unconditionality and polar duality give
\[
 \|e\|_K=\sup_{\lambda\in P_+}\lambda^Te.
\]
Indeed the polar is unconditional, and replacing a polar vector by
its coordinatewise absolute value can only increase its pairing with
$e\ge0$. Select feasible polar directions $\lambda_s$ whose images
form a spanner with coefficient bound $c$, and put
$\Psi=\sum_s\lambda_s^TF$. All these scalar functions are convex.
The same signed-coefficient argument as in \cref{lem:curvature-spanner}
gives, on every interval,
\begin{equation}\label{eq:polar-gap-domination}
 \|g_F\|_K\le c g_\Psi.
\end{equation}
The selected sum is nonaffine when $r>0$: if its nonnegative summand
gaps vanished, each selected scalarization would be affine, contradicting
independence in the nonlinear image.

Let $B$ have independent feasible columns in $C$ giving coefficient
bound $d$. Symmetry and convexity imply
\[
 B\{u:\|u\|_1\le1\}\subset C\subset B[-d,d]^r.
\]
Put $P_0=r^{-1}B[-1,1]^r$ and $P=VP_0$. Then
\begin{equation}\label{eq:inner-image-band}
 P\subset K\cap S\subset drP,
 \qquad w-y\in P/2
 \ \Longleftrightarrow\ w-y=VBu/(2r),\quad -1\le u_i\le1.
\end{equation}
The right-hand band is linear with continuous $u$.

For the finite theorem take maximum-absolute-determinant bases in the
compact spanning positive-polar image and in $C$. Cramer's rule gives
$c=d=1$. On every parity hull, each selected polar function has midpoint
gap at most one, so $g_\Psi\le2r$. Refine to scalar tolerance
$1/(2r)$ with at most $8r^2-1$ cells. Equations
\eqref{eq:polar-gap-domination} and \eqref{eq:inner-image-band} put
the chord gap in $P/2$. The exact chord center with band $P/2$ contains
the graph and admits errors in $P\subset K$. Finite disjunction proves
\eqref{eq:oracle-rank-count}. Only the midpoint Jensen vectors, which
automatically lie in $S$, are used from the comparator lift; its other
admitted errors may lie outside $S$.

For the rational theorem use the two factor-three spanners just
constructed, so $c=d=3$. Set $\tau=1/(36r)$. Refinement gives
$N_\tau(\Psi)\le(144r^2-1)2^p$. A hybrid cell has
$g_\Psi\le13/(576r)$, hence its vector gap belongs to
\[
 (13/(192r))K\cap S\subset(13/64)P.
\]
Evaluate the effective coordinate polynomials $q$ at the common knots
and round each coordinate within
$\delta=(16rL_B)^{-1}$, where
$L_B=1+\sum_{ij}|(B^{-1})_{ij}|$.
Each rounding vector belongs to $P_0/16$, because its image under
$B^{-1}$ has infinity norm at most $1/(16r)$. Interpolation preserves
this property. Restore the affine part exactly to get
$y=a+bx+Vy_q$. Then $y-F\in(17/64)P$.
The explicit band in \eqref{eq:inner-image-band} contains the graph
and admits errors in $(49/64)P\subset K$.
Rounding in effective coordinates keeps every error exactly in $S$.
The count is less than
$486\cdot144r^2 2^p=69984r^2 2^p<2^{17}r^2 2^p$.
The fixed-grid spanner bounds, rational coordinate evaluations, and
\cref{lem:indexed-compiler} prove the encoding and time claims.
\end{proof}

\section{Separable coupled vector outputs}\label{sec:separable-vector}

Let
\begin{equation}\label{eq:separable-vector-model}
 F_j(x)=\ell_j(x)+\sum_{i=1}^n\phi_{ji}(x_i),
 \qquad x\in[0,1]^n,
\end{equation}
where every $\phi_{ji}$ is continuous and convex, and every $\ell_j$
is affine. Remove coordinates on which all outputs are affine; $n$
below counts the remaining coordinates. In the computational assertions
the displayed separated representation is supplied, with rational affine
terms and dense rational polynomial summands. Mixed-coordinate nonlinear
terms are not covered by this representation.

For box tolerances, form one row for each normalized output by concatenating
its classes $\phi_{ji}/\eps_j$ modulo affine functions over all input
coordinates. Let $r$ be the rank of these rows. For explicit facets use
the concatenated rows of $(AF)_k/b_k$ instead. For a general unconditional
body, subtract the endpoint chord from each vector summand and let $S$
be the span of all resulting vectors over all coordinates; set $r=\dim S$.
For polynomial data these are ordinary rational coefficient-matrix ranks.
In particular, the basis is shared across coordinates: selecting unrelated
output bases for different inputs would not justify the comparison below.

\begin{lemma}[Product packing for a shared scalarization]\label{lem:vector-product-packing}
Suppose $\Psi=\text{affine}+\sum_i\psi_i(x_i)$ is a sum of $r$ original
normalized outputs or feasible positive-polar scalarizations. For a local
tolerance $\tau>0$, let $P_i$ be the size of a maximal scalar packing
with pairwise $J_{\psi_i}>\tau$. If $q=r/(n\tau)$ is a positive integer,
then
\begin{equation}\label{eq:vector-product-packing}
 2^{\pconv}\ge\frac{\prod_iP_i}{[3(2q+1)]^n}.
\end{equation}
A finite chord construction has local binary capacities at most $12P_i$;
a dense hybrid construction has capacities at most $5832P_i$.
\end{lemma}
\begin{proof}
Order each scalar packing. By \cref{lem:jensen-superadditivity},
an index difference $h>0$ gives midpoint gap greater than $h\tau$.
Two distinct product points with index vectors $u,v$ therefore have
$J_\Psi>\tau\|u-v\|_1$. If this exceeds $r$, one selected normalized
output or polar function has midpoint gap greater than one, making the
original vector graph midpoint inadmissible.
Greedily retain product indices and delete the integer $\ell_1$ ball
of radius $qn$ after each choice. With $t=q/(q+1)$, its size satisfies
\[
 \#\{z\in\Z^n:\|z\|_1\le qn\}
 \le t^{-qn}\left(\sum_{k\in\Z}t^{|k|}\right)^n
 =(1+1/q)^{qn}(2q+1)^n<[3(2q+1)]^n.
\]
The retained set has at least the product cardinality divided by this
bound, and pairwise inadmissible midpoints. Parity proves
\eqref{eq:vector-product-packing}; neither the packings nor this code
are computational inputs. Finally \cref{lem:scalar-chords} gives
$N_\tau(\psi_i)\le6P_i$. Finite binary encoding has capacity
$2^{\lceil\log_2N_\tau\rceil}\le2N_\tau\le12P_i$.
The hybrid has at most $486N_\tau$ cells and capacity at most
$972N_\tau\le5832P_i$, including one-cell cases.
\end{proof}

\begin{theorem}[Separable vector rank comparisons]\label{thm:separable-vector-rank}
For $n,r\ge1$ in \eqref{eq:separable-vector-model}, the finite and
polynomial rational bounds are as follows:
\begin{center}
\begin{tabular}{lll}
\toprule
Error description & Finite upper for $\pbin-\pconv$ & Compiled upper for $\pout-\pconv$\\
\midrule
Box & $n\lceil\log_2r\rceil+7n$ & $n\lceil\log_2r\rceil+17n$\\
Nonnegative facets & $n\lceil\log_2r\rceil+8n$ & $n\lceil\log_2r\rceil+18n$\\
Unconditional oracle & $2n\lceil\log_2r\rceil+8n$ & $2n\lceil\log_2r\rceil+21n$\\
\bottomrule
\end{tabular}
\end{center}
The finite oracle-body statement requires only the body itself; the
compiled statement uses the strong rational oracle and known rational
inner/outer radii of \cref{thm:vector-oracle-rank}. Rank zero or no
active input coordinates means an exact affine graph with no integers.
\end{theorem}
\begin{proof}
For a box select one exact or factor-two barycentric spanner of the
concatenated original normalized output rows. Its representation
coefficients apply in every coordinate. Summing the selected functions
gives $\Psi=\text{affine}+\sum_i\psi_i$, with
$0\le g_{ji}\le c g_{\psi_i}$, where $c=1$ in the finite case and
$c=2$ in the compiled case. If some $\psi_i$ is affine, every
corresponding original summand is affine, so that coordinate was removable.

For the finite box construction use $\tau=1/n$ and hence $q=r$ in
\cref{lem:vector-product-packing}. The packing denominator is at most
$(9r)^n$. Choose finite coordinate partitions of error $\tau$ and
interpolate every original summand at the same coordinate knots and
weight. Each normalized output's total gap is between zero and one.
The band $T_j-1\le w_j\le T_j$ contains the graph and admits unit
normalized error. One finite disjunction per coordinate encodes that
coordinate's input and all its component chords. Adding the output
bands and restoring affine terms adds only continuous equations.
The count is at most $\pconv+n\log_2(12\cdot9r)$, which is less
than $\pconv+n\log_2r+7n$.

For the compiled box construction set $\tau=1/(2n)$, so $q=2r$
and the packing denominator is at most $(15r)^n$.
Each hybrid coordinate cell has $g_{\psi_i}\le13/(32n)$, implying
$0\le\sum_i g_{ji}\le13/16$. Round each normalized original
summand endpoint downward within $1/(8n)$ and sum interpolants to
obtain $y_j$, restoring the affine part exactly. The total downward
rounding is at most $1/8$. The band
$y_j-13/16\le w_j\le y_j+1/8$ contains the graph and admits error
at most $15/16$. The count factor per coordinate is
$5832\cdot15r=87480r<2^{17}r$.

For explicit facets apply the same common-row construction to the
normalized facet functions. In the finite case choose $\tau=1/(2n)$,
so the total gap of each facet is at most $1/2$. Thus $g_F\in K/2$,
and the band $w-T\in K/2$ works. The packing denominator is
$(15r)^n$ and the count factor is $12\cdot15r<2^8r$.
In the compiled case use $\tau=1/(4n)$, so the packing denominator
is at most $(27r)^n$. Summing factor-two gap bounds gives
$g_F\in(13/32)K$. Round each original summand endpoint within
$\min\{1,(16nM_A)^{-1}\}$, with $M_A$ defined in the proof of
\cref{thm:vector-rank}. The total rounding error belongs to $K/16$;
therefore the center differs from $F$ by an element of $(15/32)K$.
Use \eqref{eq:facet-linear-band}. The count factor is
$5832\cdot27r=157464r<2^{18}r$.

For an arbitrary unconditional body choose the same effective coordinates
and two spanners as in \cref{thm:vector-oracle-rank}, now for the shared
image of all coordinate blocks. Write
$F=\ell+V\sum_iq_i(x_i)$ and
$\Psi=\sum_s\lambda_s^TF=\text{affine}+\sum_i\psi_i$.
Equations \eqref{eq:polar-gap-domination} and
\eqref{eq:inner-image-band} hold for every coordinate gap, using the
same $c,d$ and $P$. In the finite case $c=d=1$ and choose
$\tau=1/(2rn)$, so $q=2r^2$. The total gap has $K$ norm at most
$1/(2r)$ and lies in $S$, hence belongs to $P/2$.
The exact chord center with band $P/2$ works. The packing denominator
is at most $(15r^2)^n$ and the count factor is
$12\cdot15r^2<2^8r^2$.

For the compiled oracle case $c=d=3$ and set $\tau=1/(36rn)$, so
$q=36r^2$. The packing denominator is at most $(219r^2)^n$.
The total scalar hybrid gap is at most $13/(576r)$, and the vector
gap belongs to $(13/64)P$. Round every entry of each $q_i$ at its
knots within $(16nrL_B)^{-1}$. Its error belongs to $P_0/(16n)$;
summing and restoring the affine part gives a center with
$y-F\in(17/64)P$. The explicit band in
\eqref{eq:inner-image-band} admits only $(49/64)P\subset K$.
The count factor is $5832\cdot219r^2=1277208r^2<2^{21}r^2$.
In every construction all outputs share each coordinate's index and
continuous interpolation weight. All extra evaluations, basis maps,
rounding data and bands have polynomial rational size in the specified
computational models. The extra factors $n,r$ in tolerances cost only
logarithmic additional precision. Compact nonnegative-facet bodies and
positive-polar domination give the asserted affine degeneracies.
\end{proof}

For weighted $\ell_1$ error the single facet has rank one whenever the
graph is nonlinear. The separable finite and compiled overheads are
therefore $8n$ and $18n$. General mixed-coordinate polynomials require
additional arguments: replacing the sum of univariate chords by
multilinear corner interpolation does not itself supply the necessary
compact linear encoding.

\section{Limits of scalarization and simultaneous refinement}
\label{sec:vector-obstructions}

The rank comparisons do not follow by taking the largest of the separate
scalarized minimum dimensions. The following family makes every such
scalar minimum small while the simultaneous integer dimension grows.

\begin{proposition}[Separated positive-power layers]\label{prop:vector-power-obstruction}
For $M\ge1$, put
\[
 D_j=128\cdot1024^{j-1},\quad F_j(x)=\tfrac74x^{D_j},
 \quad x_j=1-64/D_j,\qquad K=[-1,1]^M.
\]
Every exact-graph midpoint has component error strictly below $7/8$.
For every real $\lambda$ with $\|\lambda\|_1=1$, the scalar graph
$\lambda^TF$ at its induced tolerance one has a finite binary linear
formulation with at most one binary. Nevertheless
\begin{equation}\label{eq:power-vector-counts}
 \lceil\log_3M\rceil\le\pconv\le\pbin,
 \qquad
 \lceil\log_2M\rceil\le\pbin\le\lceil\log_2(M+1)\rceil.
\end{equation}
If $N_F(\eta)$ is the least number of interval chords with vector error
at most $\eta$ in the infinity norm, then $N_F(2)=1$ and $N_F(1)\ge M$.
The binary lower bound holds even for arbitrary convex binary lifts.
\end{proposition}
\begin{proof}
For convex differentiable $f$, $J_f(a,b)$ is nonincreasing in $a$
and nondecreasing in $b$, as follows by differentiating the midpoint
formula. Hence $0\le J_{F_j}(a,b)\le J_{F_j}(0,1)
=7/8-(7/4)2^{-D_j}<7/8$.

For $j<k$, Bernoulli's inequality gives
$x_k^{D_j}\ge1-64D_j/D_k\ge15/16$. At
$\xi=(x_j+2x_k)/3\le1-64/(3D_j)$,
$\xi^{D_j}\le e^{-64/3}\le3/67$. Therefore the thirds chord error
in component $j$ is at least
\begin{equation}\label{eq:power-thirds-violation}
 \frac74\left(\frac23\frac{15}{16}-\frac3{67}\right)
 =\frac{2177}{2144}>1.
\end{equation}
Selected exact graph witnesses at these $M$ inputs must have distinct
integer residues modulo three; otherwise their indicated weighted
average has integer coordinates and violates the error body.
Thus $3^p\ge M$. In a binary lift, equal assignments admit every
convex weight, so the same witnesses require $2^p\ge M$.
These arguments use neither bounded integer labels nor closedness.

Each component increases from zero to $7/4$. Its half-height point
$a_j=2^{-1/D_j}$ splits its range into intervals of length $7/8$.
Overlay all half-height points to obtain $M+1$ input intervals.
On each, the rectangle formed from the input interval and all component
ranges contains the vector graph and has component error at most $7/8$.
Finite disjunction proves the upper binary count; its coefficients may
be real algebraic. The full-domain chord error is at most $7/4<2$.
An interval with unit vector chord error cannot contain two selected
$x_j,x_k$: its component chord lies above the chord through those two
contacts and would inherit \eqref{eq:power-thirds-violation}.
Thus $N_F(1)\ge M$.

Finally set $\Psi_\lambda=\sum_j|\lambda_j|F_j$. It increases
continuously from zero to $7/4$, and
$|\lambda^TF(b)-\lambda^TF(a)|\le\Psi_\lambda(b)-\Psi_\lambda(a)$
for $a\le b$. Split at $\Psi_\lambda=7/8$ and use the scalar
input-output range rectangle on each side. Each range has length at
most $7/8$, giving one binary. Negative scalarization coefficients
are allowed. The induced scalar tolerance is exactly
$h_K(\lambda)=\|\lambda\|_1=1$.
\end{proof}

The statement quantifies separately over scalarizations: each may choose
its own partition and formulation. It does not question polar duality,
which exactly describes the norm, or rule out algorithms combining
several scalarizations. It also does not establish an unbounded difference
$\pbin-\pconv$. The family has curvature rank $M$, so its behavior is
consistent with the rank-dependent upper bounds.

\begin{proposition}[Cap-set contact restriction]\label{prop:cap-set}
For the family of \cref{prop:vector-power-obstruction}, let
$\operatorname{cap}(p)$ be the largest subset of $\mathbb F_3^p$ with
no three distinct elements summing to zero. Every $p$-integer convex
lift satisfies
\begin{equation}\label{eq:cap-count}
 M\le\operatorname{cap}(p)\le3c_*^p,
 \quad c_*=(1+t_*+t_*^2)t_*^{-2/3}<3,
 \quad t_*=(\sqrt{33}-1)/8.
\end{equation}
In particular $\pconv\ge\max\{0,\lceil\log_{c_*}(M/3)\rceil\}$,
and for $M=3$ the exact counts are $\pconv=\pbin=2$.
\end{proposition}
\begin{proof}
The selected witness residues are distinct by
\eqref{eq:power-thirds-violation}. If three distinct residues summed
to zero, their uniform lifted average would be integral. Let $j$ be
the smallest selected input index. Its weight is $1/3$ and the other
two inputs have component-$j$ powers at least $15/16$. The average
input is at most $1-64/(3D_j)$. Equation
\eqref{eq:power-thirds-violation} again gives a forbidden error.
Thus the residues form a cap set. Over $\mathbb F_3$, a zero-sum
triple with two equal entries has all three equal, so this is precisely
the absence of nontrivial three-term progressions.

Theorem 4 of \citet{ellenberg2017}, specialized to $\mathbb F_3$,
bounds its size by
$3\#\{\alpha\in\{0,1,2\}^p:\sum_i\alpha_i\le2p/3\}$.
For $0<t<1$, this monomial count is at most
$t^{-2p/3}(1+t+t^2)^p$, by bounding the indicator with
$t^{\sum\alpha_i-2p/3}$ and summing. Differentiating its base gives
$4t^2+t-2=0$ at the unique minimum, proving \eqref{eq:cap-count}.
For a simpler rational base, $t=1/2$ gives
$(7/4)2^{2/3}<14/5$, so $M\le3(14/5)^p$.
A cap set in $\mathbb F_3$ has at most two elements. Thus $M=3$
needs at least two general integers, while
\eqref{eq:power-vector-counts} gives two binaries.
\end{proof}

There is also a direct one-integer obstruction which illustrates the role
of repeated convexification. For any selected $x_j<x_k$, every chord
weight $t\in[2/3,4/5]$ on the larger input has component-$j$ gap at least
\[
 \frac74\left(\frac23\frac{15}{16}-\frac{25}{2048}\right)
 =\frac{8785}{8192}>1,
\]
because the averaged input is at most $1-64/(5D_j)$ and
$e^{-64/5}\le25/2048$, using $e^v\ge v^2/2$.
If two integer labels differ in magnitude by $q\ge3$, some $k/q$
lies in this weight interval: choose $k=\lceil2q/3\rceil$;
$q=3,4$ are direct and $q\ge5$ gives
$k\le(2q+2)/3\le4q/5$. The corresponding chord has an integer
label and is forbidden. Equal labels are forbidden as well, so three
selected labels would have to be three consecutive integers. The midpoint
of the extreme labeled witnesses belongs to the middle integer section.
Mix it with the third witness so that each original witness has weight
$1/3$. This stays in that section and violates
\eqref{eq:power-thirds-violation}. Checking only integer points on
segments between initially assigned labels would miss this second
convexification. The cap-set theorem supplies the higher-dimensional
version; its exponential base greater than two and the absence of a
matching upper construction leave the binary/general-integer gap unresolved.

\section{Exact separations and the role of error geometry}
\label{sec:exact-separations}

\subsection{Fixed-degree convex graphs with exact product counts}

\begin{theorem}[Exact convex box counts]\label{thm:exact-box-gap}
For $n\ge1$, let
\[
 F_n(x)=\left(\tfrac74(1-x_i)^{32},\tfrac74x_i^{32}\right)_{i=1}^n,
 \qquad x\in[0,1]^n,
\]
with unit error in each of the $2n$ components. Then
\begin{equation}\label{eq:exact-box-counts}
 \pconv(F_n)=n,\qquad \pbin(F_n)=\lceil n\log_2 3\rceil.
\end{equation}
The lower binary count is valid for arbitrary convex binary lifts.
Requiring closed convex lifts leaves both minima unchanged.
Both upper bounds have rational MILP realizations; the general-integer
one uses $3n$ continuous auxiliaries and $13n$ inequalities, with rational
coefficients from a fixed finite set independent of $n$. The binary
construction asserted here may have exponential continuous size.
The linear size claim concerns variable and row counts, without a claim
of linear serialized bit length or a solver running-time bound.
\end{theorem}
\begin{proof}
For one input put $A=7/4$, $c=1/48$, $L=A(1-c)^{32}$,
$d=Ac^{32}$, and $Q(x)=(A(1-x)^{32},Ax^{32})$.
We have $A-1<L<1$, $d<1/4$, and $(A+d)/2<1$.
Indeed Bernoulli gives $(47/48)^{16}\ge2/3$ and hence
$L\ge7/9>3/4$. The binomial expansion gives
$(48/47)^{32}\ge1+32/47+496/47^2>7/4$, so $L<1$.
The remaining inequalities follow from $d<A/48<1/4$.
Consider the rational boxes
\begin{align*}
 B_0&=[0,c]\times[L,A]\times[0,d],\\
 B_1&=[c,1-c]\times[0,L]^2,\\
 B_2&=[1-c,1]\times[0,d]\times[L,A].
\end{align*}
Each contains the graph on its input interval and has unit graph error:
the outer component range widths are $A-L<1$ and $d<1$, and both
true and admitted middle outputs lie in $[0,L]$.
Append labels $0,1,2$ and take their convex hull $\mathcal C$.
At integer label zero or two only the corresponding box remains.
At label one, consolidating within each box gives weights
$t,1-2t,t$, $0\le t\le1/2$. Its input is at least
$(1-2t)c+t(1-c)=c+t(1-3c)\ge c$ and, by symmetry, at most $1-c$.
Every output is in
$[0,(1-2t)L+t(A+d)]\subset[0,\max\{L,(A+d)/2\}]\subset[0,1)$.
Every true output there also lies in $[0,L]$. Thus the entire middle
integer section, including all mixtures, is valid.

Write $B_j=[\ell_j,u_j]$ coordinatewise in $v=(x,w_1,w_2)$.
An explicit rational linear lift of $\mathcal C$ is
\begin{align*}
 &\lambda_j\ge0\quad(j=0,1,2),\qquad
   \sum_{j=0}^2\lambda_j=1,\qquad z=\lambda_1+2\lambda_2,\\
 &\sum_{j=0}^2\lambda_j\ell_{j,k}\le v_k
   \le\sum_{j=0}^2\lambda_j u_{j,k}\quad(k=1,2,3).
\end{align*}
For fixed weights, the weighted sum of boxes is exactly the box with
these weighted endpoints, since its coordinate intervals vary independently.
Hence projection gives precisely the labeled hull. Only $z$ is integer.
There are three continuous auxiliaries, nine inequalities and two equalities,
or thirteen inequalities after splitting the equalities. Products prove
$\pconv(F_n)\le n$ with the stated counts and fixed rational data.

The three contacts at $0,1/2,1$ are pairwise incompatible under an
oriented thirds combination. For $0,1/2$, give weight $2/3$ to zero;
the first-component gap at $1/6$ is
\[
 \tfrac23A+\tfrac13A2^{-32}-A(5/6)^{32}>1.
\]
Since $4\cdot5^8<6^8$, $(5/6)^{32}<1/256$, and the gap exceeds
$7/6-7/1024>1$. The pair $1/2,1$ is symmetric. For $0,1$, weight
$2/3$ on zero gives gap $\tfrac23A-A(2/3)^{32}>1$ by the same bound.

Choose witnesses at all $3^n$ product inputs in $\{0,1/2,1\}^n$.
Two integer codes equal modulo three would permit both thirds
combinations. They differ in some input coordinate, where one orientation
is forbidden by the preceding calculation. Thus all codes have distinct
residues and $3^p\ge3^n$, proving $p\ge n$.
Equal binary codes would permit every convex weight, so these same
contacts need distinct binary assignments: $2^p\ge3^n$.
For equality, take the $3^n$ products of the three valid boxes and
apply \cref{lem:disjunction} with $\lceil\log_2(3^n)\rceil$ bits.
This is a finite rational polyhedral construction.
Both rational upper constructions use closed lifted sets, while the lower
bounds allow all convex lifts, proving the closedness assertion as well.
\end{proof}

Adding $56x_i$ to the first component of each pair makes both components
nondecreasing: its derivative becomes $56[1-(1-x_i)^{31}]\ge0$.
This affine shear preserves errors, both minima, and the closed-lift
minima. Substitution in the displayed formulation preserves its $3n$
auxiliary and $13n$ inequality counts. It does not make all monomial
coefficients nonnegative: the first component's cubic coefficient is $-8680$.
Equation \eqref{eq:exact-box-counts}
shows that an additive loss proportional to input dimension can be
necessary even with fixed degree, tolerance and numerical data. It does
not amplify the gap with only one input.

A piecewise-linear precursor explains the three-section geometry.
Take $a=1/8$, $A=3/2$ and
$H(x)=(A\max\{1-x/a,0\},A\max\{(x-1+a)/a,0\})$.
Its three exact affine graph segments, labeled $0,1,2$, have a valid
one-integer convex hull at unit box error. In the middle integer section,
the outer weights are $t,t$ and their averaged input lies in
$[(1-a)/2,(1+a)/2]\subset[a,1-a]$. The true graph there is zero;
the admitted outputs are between zero and $tA\le3/4$.
The contacts $0,1/2,1$ cannot share a binary section pairwise: the
first pair's chord at $a$ has first output $A(1-2a)=9/8$, the last
pair is symmetric, and the outer-pair chord gives $A(1-a)=21/16$.
Thus this hinge curve also has exact counts one and two.

This strict separation survives a convexity-preserving polynomial
approximation. For a scalar $L$-Lipschitz function $h$, its Bernstein
polynomial $B_Nh(x)=\sum_kh(k/N)\binom Nkx^k(1-x)^{N-k}$ satisfies
\begin{equation}\label{eq:bernstein-lipschitz}
 |B_Nh(x)-h(x)|
 \le L\,\mathbb E|X/N-x|
 \le L\sqrt{x(1-x)/N}\le L/(2\sqrt N),
 \quad X\sim\operatorname{Bin}(N,x).
\end{equation}
For convex $h$, its sampled second differences are nonnegative, and
\[
 (B_Nh)''=N(N-1)\sum_{k=0}^{N-2}
 [h((k+2)/N)-2h((k+1)/N)+h(k/N)]
 \binom{N-2}{k}x^k(1-x)^{N-2-k}\ge0.
\]
The hinge components are $12$-Lipschitz. Taking $N=2^{18}$ gives
uniform error at most $3/256<\delta=1/64$. Thicken the labeled
polytope in outputs by $[-\delta,\delta]^2$. It contains the exact
Bernstein graph and its integer-section error is at most
$3/4+2\delta<1$. Each forbidden chord loses at most $2\delta$,
so the binary lower remains two; the three separately thickened segments
give two binaries. This is a distinct stability mechanism, although the
degree-32 family in \cref{thm:exact-box-gap} is simpler and sharper.
The probabilistic Bernstein estimate and its convexity-preserving
second-difference identity are classical approximation tools.

\subsection{Nonconvex scalar polynomials and degree dependence}

\begin{theorem}[Two general integers and growing binary dimension]
\label{thm:nonconvex-gap}
For every integer $M\ge2$ there is a dense rational polynomial $q_M$
of degree at most $1024M^2$ such that, at fixed absolute error $1/4$,
\begin{equation}\label{eq:nonconvex-gap-counts}
 \pconv(q_M;1/4)\le2,\qquad
 \lceil\log_2M\rceil\le\pbin(q_M;1/4)
 \le\lceil\log_2M\rceil+1.
\end{equation}
The lower bound holds even for arbitrary convex binary lifts.
For a continuous scalar function with $s$ convex or concave pieces,
\begin{equation}\label{eq:convexity-piece-bound}
 \pbin\le\pconv+\lceil\log_2(3s)\rceil.
\end{equation}
Consequently a degree-$D$ nonaffine polynomial has finite overhead
$\lceil\log_2(3\max\{1,D-1\})\rceil$.
\end{theorem}
\begin{proof}
Let $T_M$ be the triangular wave of $M$ periods: for
$t=Mx-j\in[0,1]$, $j=0,\ldots,M-1$, set
$T_M(x)=2t$ for $t\le1/2$ and $T_M(x)=2-2t$ otherwise.
It is $2M$-Lipschitz. Take $N=(32M)^2$ and $q_M=B_NT_M$.
Equation \eqref{eq:bernstein-lipschitz} gives
$|q_M-T_M|\le1/32$, and positivity of Bernstein weights gives
$0\le q_M\le1$. Sampled values are rational of polynomial encoding;
binomial expansion gives a dense monomial representation of polynomial
bit length in $N$ and $\log M$.

An integer period $z\in\{0,\ldots,M-1\}$ and one binary orientation
$b$ encode the triangular-wave graph exactly by $x=(z+t)/M$ and the
two bounded linear branches
\[
 b=0:\ 0\le t\le1/2,\ y=2t;
 \qquad b=1:\ 1/2\le t\le1,\ y=2-2t.
\]
A two-branch disjunction with finite rational big-$M$ bounds realizes
these equations. Period boundaries and $x=1$ are included. Adding
$|w-y|\le1/32$ contains the entire $q_M$ graph and admits error
at most $1/16<1/4$. The period and orientation together count as two
general integers.

At the $M$ peaks $(j+1/2)/M$, $q_M\ge31/32$; at each intervening
trough $k/M$, $q_M\le1/32$. If two peak contacts had witnesses in
the same binary section, their full chord would be admitted. At a
trough between them its value is at least $31/32$, yielding error
at least $30/32>1/4$. Thus peaks need distinct binary codes.
Conversely replace the period integer in the explicit upper construction
by $\lceil\log_2M\rceil$ digits and exclude unused codes, retaining
the orientation bit. This proves \eqref{eq:nonconvex-gap-counts}.

For \eqref{eq:convexity-piece-bound}, restrict any $p$-integer lift
to each shape interval and negate a concave output. The scalar chord
lemma gives at most $3\cdot2^p$ error-valid chord bands per piece.
Their finite union has at most $3s2^p$ members. A polynomial of
degree $D\ge2$ has at most $D-2$ distinct roots of its second
derivative, and hence at most $D-1$ such intervals. An affine
polynomial is exact with zero integers. Roots may be irrational in this
finite real-coefficient argument.
\end{proof}

The polynomial degrees in this family tend to infinity: alternating
peaks and troughs force at least $2M$ distinct zeros of $q_M-1/2$.
Together with the upper degree bound, \eqref{eq:nonconvex-gap-counts}
shows that logarithmic degree overhead is the correct worst-case order.
The general binary/general-integer distinction and binarization of
bounded integer variables are established; the example imposes the
additional requirement of one connected rational polynomial graph at
fixed positive tolerance. It is a lower bound on optimum counts, not
only on one approximation construction.

\begin{theorem}[Compiled arbitrary polynomial vectors]\label{thm:arbitrary-vector-overlay}
Let the nonaffine outputs of $F:[0,1]\to\R^m$ be dense rational
polynomials of degrees $D_j\ge2$, with positive rational box tolerances.
A polynomial-time rational MILP construction has
\begin{equation}\label{eq:arbitrary-vector-count}
 \pout\le\pconv+12+\left\lceil\log_2\sum_{j:D_j\ge2}D_j\right\rceil.
\end{equation}
Affine outputs are represented exactly, and an entirely affine graph
uses no integers. For one output this gives overhead
$12+\lceil\log_2D\rceil$.
\end{theorem}
\begin{proof}
For output $f_j=\sum_kc_{jk}x^k$ put
$M_j=\max\{1,\sum_{k\ge1}k|c_{jk}|\}\ge\max|f_j'|$.
Isolate each distinct interior zero of $f_j''$ in a disjoint rational
bracket of width at most $\min\{1,\eps_j/(32M_j)\}$, with nonroot
endpoints. Squarefree preprocessing, root isolation and refinement have
polynomial bit cost by the univariate root routine used in
\cref{thm:convex-hybrid}; endpoint roots need no bracket.
The brackets and complementary intervals give at most $2D_j$ pieces.
On a complement the curvature has constant sign; normalize its rational
interval and apply the convex hybrid to $f_j$ or $-f_j$.
Restriction and output negation preserve the comparator's integer count.
Each such piece has fewer than $2^{p+11}$ cells, where $p$ is the
integer dimension of any comparator vector lift. A bracket has one cell:
a derivative bound gives absolute chord error at most
$2M_jh\le\eps_j/16$, including on every subinterval.
Thus the scalar list has
\begin{equation}\label{eq:arbitrary-scalar-cells}
 K_j\le4096D_j2^p.
\end{equation}

Use ordered local quantiles, reverse reflected indices, and concatenate
in original input order. The indexed evaluator records source-cell type:
convex, concave, or bracket. All knots have a polynomial-bit common
denominator as in \cref{lem:implicit-overlay}. Merge their right
endpoints with multiplicity. On every positive merged cell recover each
output's containing source cell and type using the search in that lemma;
at a repeated endpoint any containing source cell is valid.

Evaluate each original polynomial at the selected merged endpoints.
Round downward within $\eps_j/16$ for convex cells and brackets,
and upward within $\eps_j/16$ for concave cells. If $y_j$ is the
interpolant, use
\begin{equation}\label{eq:signed-overlay-bands}
 \begin{cases}
 y_j-13\eps_j/16\le w_j\le y_j+\eps_j/8,
       &\text{convex or bracket},\\
 y_j-\eps_j/8\le w_j\le y_j+13\eps_j/16,
       &\text{concave}.
 \end{cases}
\end{equation}
For convex cells $y_j-f_j\in[-\eps_j/16,13\eps_j/16]$;
for concave cells it lies in $[-13\eps_j/16,\eps_j/16]$.
For brackets it lies in $[-\eps_j/8,\eps_j/16]$.
Each band contains the graph and has absolute admitted error at most
$15\eps_j/16$. The same checks hold on singleton cells. The type
flags are computed Boolean wires and choose the offsets linearly;
they require no new integer declarations. A common interpolation weight
and input cell guarantee simultaneous vector graph containment.

The total merged count is $K=\sum_jK_j$, so
\eqref{eq:arbitrary-scalar-cells} proves
\eqref{eq:arbitrary-vector-count}. All numerator searches, source-cell
lookups, directed evaluations, flags and invalid-code exclusions have
polynomial bit complexity. \Cref{lem:indexed-compiler} declares only
the $\lceil\log_2K\rceil$ input index bits integer. The algorithm
uses actual computed counts, not $p$ or the comparator lift.
\end{proof}

\subsection{Tilted error bodies and conditioning}

\begin{proposition}[A two-output tilted-body gap]\label{prop:tilted-gap}
There are rational componentwise convex polynomial curves with two outputs
and rational symmetric parallelogram error bodies for which
$\pconv\le2$ and $\pbin\ge\lceil\log_2M\rceil\to\infty$.
The Euclidean outer-to-inner radius ratio of these bodies is at least
$8+60M^2$; this is not a uniformly conditioned family.
\end{proposition}
\begin{proof}
Use $q_M$ from \cref{thm:nonconvex-gap}, with dense coefficients $c_k$,
and put
\[
 C_M=1+\sum_{k\ge2}k(k-1)|c_k|,\quad
 F_M=(q_M+C_Mx^2,C_Mx^2),\quad
 K_M=\{e:|e_1-e_2|\le1/4,\ |e_2|\le C_M\}.
\]
Since $|q_M''|\le C_M-1$, both components are convex. The body is
compact, symmetric and has interior, but is not unconditional: it
contains $(C_M,C_M)$ and not its first-coordinate sign flip.
Take the two-integer scalar lift in \cref{thm:nonconvex-gap}, with
output $s$ of error at most $1/16$. Add $0\le w_2\le C_M$ and
$w_1=s+w_2$. Exact vector graph points lift with
$s=q_M(x)$, $w_2=C_Mx^2$; every admitted error obeys
$|e_1-e_2|\le1/16$ and $|e_2|\le C_M$. This gives the general-integer
upper bound with rational linear constraints.
Any binary convex lift projects by $s=w_1-w_2$ to a scalar lift of
$q_M$ with tolerance $1/4$. Its binary count is at least
$\lceil\log_2M\rceil$; replacing the period integer also gives an
upper count of $\lceil\log_2M\rceil+1$.

The body has circumradius at least $\sqrt2C_M$ and inradius at most
$1/(4\sqrt2)$ because of its narrow difference strip. At a peak $a$
and adjacent troughs $a\pm h$, $h=1/(2M)$, the uniform Bernstein
estimate gives
$q_M(a-h)-2q_M(a)+q_M(a+h)\le-15/8$.
The identity
\[
 q_M(a-h)-2q_M(a)+q_M(a+h)
 =\int_{-h}^h(h-|t|)q_M''(a+t)\,dt
\]
implies $\max|q_M''|\ge(15/2)M^2$. Thus
$C_M\ge1+(15/2)M^2$, proving the claimed ratio.
\end{proof}

Adding a sufficiently large quadratic to make a smooth function convex
is the elementary difference-of-convex construction; the general framework
predates this application \citep{hartman1959}. The large permitted error
in the common output direction absorbs that quadratic. The narrow
difference direction retains the nonconvex scalar obstruction. The error
geometry, as well as componentwise convexity, is therefore material.

\begin{proposition}[Finite conditioning comparisons]\label{prop:conditioning}
Let $F:[0,1]\to\R^m$ be continuous and componentwise convex.
If $aB_\infty^m\subset K\subset bB_\infty^m$, $0<a\le b$, then
\begin{equation}\label{eq:infinity-conditioning}
 \pbin\le\pconv+
 \left\lceil\log_2\left(2\left\lceil2mb/a\right\rceil-1\right)\right\rceil.
\end{equation}
The same bound holds with Euclidean inner and outer balls. In that
Euclidean case there is also the bound
\begin{equation}\label{eq:simplex-conditioning}
 \pbin\le\pconv+
 \left\lceil\log_2\left(2\left\lceil2\sqrt{2m}\,b/a\right\rceil-1\right)\right\rceil.
\end{equation}
These are finite real-coefficient statements. They hold even for a
nonsymmetric compact convex permitted-error set containing the given
inner ball, using the same graph-sandwich definition.
\end{proposition}
\begin{proof}
On a parity hull the nonnegative midpoint vector belongs to $K$.
For $h=\sum_jF_j$, its midpoint gap is at most $mb$ under infinity
containment and at most $\sqrt m\,b$ under Euclidean containment.
The full gap is at most twice that. In the infinity case refine to
tolerance $a$, so \cref{lem:level-refinement} uses at most
$2\lceil2mb/a\rceil-1$ cells per parity hull. Each component gap
is between zero and $a$, and downward component bands have error in
$aB_\infty^m\subset K$. In the Euclidean case refine instead to
$a/\sqrt m$. Downward bands then have error in $aB_2^m$, with the
same refinement factor. Finite disjunction proves
\eqref{eq:infinity-conditioning} in both cases.

For the sharper Euclidean variant, if the total scalar gap is at most
$s$, use the simplex band
\begin{equation}\label{eq:simplex-band}
 w=T-v,\qquad v\ge0,\qquad \sum_jv_j\le s.
\end{equation}
It contains the graph by $v=g_F$. Every admitted error is $g_F-v$,
where both nonnegative vectors have $\ell_1$ norm at most $s$.
Thus $\|g_F-v\|_2^2\le\|g_F\|_2^2+\|v\|_2^2\le2s^2$.
Choose $s=a/\sqrt2$ and refine from $2\sqrt m\,b$ to $s$.
This proves \eqref{eq:simplex-conditioning}. Only the outer bound
on midpoint vectors and the inner bound for final errors were used;
symmetry of the intervening permitted-error set is unnecessary.
\end{proof}

A related normalization explains why unconditionality avoids arbitrary
conditioning losses. Let $a_j=\max_{e\in K}|e_j|>0$ and scale each
output by $1/a_j$. The normalized body $K'$ satisfies
\[
 B_1^m\subset K'\subset B_\infty^m.
\]
For the first inclusion, take a point attaining a coordinate maximum and
average its sign flips in all other coordinates. The resulting positive
and negative unit coordinate vectors lie in $K'$; their convex hull
is $B_1^m$. Consequently $m^{-1}B_\infty^m\subset K'$ and
\eqref{eq:infinity-conditioning} gives finite overhead
$\lceil\log_2(4m^2-1)\rceil$. The simplex band improves this elementary
normalization bound: now the parity-hull gap of $\sum_jF_j/a_j$ is at
most $2m$; refine to $s=1/2$. Since
$\|g_F-v\|_1\le2s=1$, the band lies in $B_1^m\subset K'$ and
has overhead $\lceil\log_2(8m-1)\rceil$.
The direct maximum-product bound in \cref{prop:vector-log-product}
is sharper, but the normalization and simplex arguments give additional
geometric comparisons without computing an optimal allocation.

\subsection{The remaining one-input box question}\label{subsec:one-input-open}

For one compact input interval and continuous componentwise convex outputs
with box error, it remains unresolved here whether a universal constant
$C$ satisfies $\pbin\le\pconv+C$, independently of output count or
curvature rank. \Cref{thm:exact-box-gap} rules out equality of the two
minima, but its unbounded product gap increases input dimension.
The positive-power family proves failure of scalar-minimum and uniform
refinement tests; neither its residue lower bound nor the cap-set
strengthening supplies a growing difference between the two minima.
The tilted-body construction changes the error geometry and loses
uniform conditioning, so it does not settle the box question either.

One can express a concrete obstacle using lattice-free convex sets.
For an admissible convex lift $\mathcal C$ and a convex component $F_j$,
the strict upper-violation set
\[
 \{(x,w,z,a)\in\mathcal C:w_j>F_j(x)+\eps_j\}
\]
is convex. Its projection to the integer coordinates is convex and has
no integer point; any such point would have an invalid feasible integer
witness. The union of these sets over outputs need not be convex.
Taking its convex hull is unjustified: mixing violations of different
components can remove every violation. Thus applying lattice-free
geometry to one output does not automatically control all outputs.

The mixed Helly identity
$h(\R^q\times\Z^p)=(q+1)2^p$ gives finite infeasibility certificates
for convex sets \citep[Theorem 1.1]{averkov2012}; it does not supply an
interval cover preserving the output error. Its dependence on the
unrestricted continuous dimension also precludes a direct count bound.
Similarly, an integral Radon partition would need a common integer index
and compatible input supports producing a violation in one output;
parity alone does not ensure those properties. These observations delimit
what the present contact arguments establish. None of the proved
comparisons depends on resolving the constant-gap question.
